\subsection{Подключение}

При неудачном подключении элементы управления остаются заблокированными. В
строке состояния отображается адрес микросервиса либо текст возникшего
исключения.

\subsection{Ручной ввод данных}

Пустая строка, неполный байт, недопустимый символ или значение вне диапазона
приводят к ошибке разбора данных. Текст ошибки отображается в строке состояния
и в поле ``Отправленные данные''.

\subsection{Выполнение SendData}

При успешной отправке строка состояния содержит сообщение
``Статус: данные отправлены''. При неуспешном результате отображается сообщение
``Статус: ошибка отправки'' и текст причины.

\subsection{Получение данных Subscribe}

При успешном получении пакета строка состояния содержит сообщение
``Статус: получены данные''. При ошибке отображаются сообщение
``Статус: ошибка приема'' и текст причины.

Если перед завершением потока была показана ошибка приема, она сохраняется в
строке состояния. В остальных случаях строка состояния сообщает:
``Статус: подписка завершена''.

\subsection{Коды ошибок микросервиса}

\begin{longtable}{|p{\dimexpr 0.35\textwidth - 12.6pt\relax}|p{\dimexpr 0.65\textwidth - 12.6pt\relax}|}
	\caption{Коды ошибок микросервиса} \\
	\hline
	\textbf{Код} & \textbf{Описание} \\
	\hline
	\endfirsthead
	\hline
	\textbf{Код} & \textbf{Описание} \\
	\hline
	\endhead
	\texttt{NO\_ERROR} & Операция выполнена успешно. \\
	\hline
	\texttt{INVALID\_CHANNEL\_ID} & Передан недопустимый идентификатор канала. \\
	\hline
	\texttt{INVALID\_DATA} & Данные отсутствуют либо имеют недопустимый формат. \\
	\hline
	\texttt{DRIVER\_NOT\_CONNECTED} & Микросервис не подключен к драйверу RS-485. \\
	\hline
	\texttt{DRIVER\_NO\_REPLY} & Драйвер или подключенное устройство не вернуло ответ. \\
	\hline
	\texttt{DRIVER\_TIMEOUT} & Истекло время ожидания ответа. \\
	\hline
	\texttt{DRIVER\_ERROR} & Драйвер сообщил об ошибке выполнения операции. \\
	\hline
	\texttt{INTERNAL\_ERROR} & Произошла внутренняя ошибка микросервиса или gRPC-взаимодействия. \\
	\hline
\end{longtable}

\subsection{Действия при ошибках}

Если элементы вкладки остаются заблокированными, необходимо:

\begin{enumerate}
	\item убедиться, что микросервис RS-485 запущен;
	\item проверить адрес сервиса в файле \texttt{rs/config/config.yaml};
	\item убедиться, что указанный TCP-порт свободен;
	\item проверить доступность конфигурационного файла для GUI;
	\item повторно запустить графический интерфейс.
\end{enumerate}

Если подключение к микросервису установлено, но отправка или подписка
завершаются ошибкой, необходимо проверить:

\begin{itemize}
	\item корректность идентификатора канала;
	\item формат вручную введенных байтов;
	\item существование и доступность выбранного файла;
	\item работу драйвера RS-485 либо \texttt{rs485\_fake\_driver};
	\item адрес драйвера в конфигурационном файле;
	\item физическое соединение RS-485 при использовании реального устройства.
\end{itemize}

При ошибках \texttt{DRIVER\_NO\_REPLY} и \texttt{DRIVER\_TIMEOUT} необходимо
дополнительно проверить доступность подключенного устройства.

При ошибке \texttt{INTERNAL\_ERROR} рекомендуется сохранить текст сообщения
из строки состояния и журналов и передать его разработчику.